\documentclass{article}
\usepackage[utf8]{inputenc}
\usepackage{amsmath, amssymb, amsthm}
\usepackage{enumitem}

\title{Appunti di Analisi T1}
\author{Leonardo}

\begin{document}

\maketitle

\section{Modulo 1}

\subsection{Principio di induzione}

Il principio di induzione è un metodo generico per dimostrare
un teorema qualsiasi con dominio in $\mathbb{N}$.

\vspace{5mm}
\noindent
Il procedimento di prova per induzione prevede un affermazione
o uguaglianza che dipende da $n$ e viene indicata con: $P(n)$.

\vspace{5mm}
\noindent
Data l'affermazione $P(n)$, il principio prova
la precedente, al verificarsi delle seguenti condizioni:

\newlist{induzione}{enumerate}{2}
\setlist[induzione]{label*=\arabic*.}

\begin{induzione}
    \item $P(0)$
    \item $P(k) \Leftrightarrow P(k+1)$
\end{induzione}

\noindent
Se queste condizioni si verificano, allora: $P(n)$ è vera $\forall n \in \mathbb{N}$.

\subsection{Principio del minimo}

\textbf{DEF}: Ogni sottoinsieme di $\mathbb{N}$ ha un minimo.

\subsection{Maggioranti e minoranti, Massimo e minimo}

Dato $\mathbb{A} = \{ a \in \mathbb{R} \}$

\newlist{maggioranti}{enumerate}{2}
\setlist[maggioranti]{label*=\arabic*.}

\begin{maggioranti}
    \item Maggiorante: $M \in \mathbb{R} \land \forall a \in \mathbb{A} : a \leq M$
    \item Minorante: $m \in \mathbb{R} \land \forall a \in \mathbb{A} : a \geq M$
\end{maggioranti}

\noindent
Se l'insieme $\mathbb{A}$ ammette un maggiorante si
dice \underline{superiormente limitato}, se ammette un minorante
si dice \underline{inferiormente limitato} e se ammette entrambe
si dice \underline{limitato}.

\vspace{5mm}
\noindent
Il minimo di un insieme $\mathbb{A}$, se esiste,
è il più grande dei \underline{minoranti}:

\vspace{5mm}
\noindent
Il massimo di un insieme $\mathbb{A}$, se esiste,
è il più piccolo dei \underline{maggioranti}:

\subsection{Piccole nozioni}

L'insieme dei numeri reali \textbf{NON} è numerabile.

\vspace{5mm}
\noindent
Radice aritmetica: $\sqrt{n} = m : m \in \mathbb{N}$.

\subsection{Successioni}

\textbf{DEF}: Una successione di numeri reali è una funzione da
$\mathbb{N}$ in $\mathbb{R}$.

\begin{equation}
    \begin{split}
        f : \mathbb{N} \longrightarrow \mathbb{R} \\
        n \longrightarrow f(n) := a_n
    \end{split}
\end{equation}

\noindent
I termini della successione sono notati nella seguente maniera:

\vspace{5mm}
$f(0) = a_0$

\vspace{5mm}
$f(1) = a_1$

\vspace{5mm}
$f(2) = a_2$

\vspace{5mm}
\noindent
La successione può essere:

\begin{itemize}
    \item Superiormente limitata
    \item Inferiormente limitata
    \item Limitata
\end{itemize}

\subsection{Intorni}

\textbf{DEF}: Si dice intorno, un insieme aperto che ammette un punto centrale.

\subsection{Limiti}

\textbf{DEF}: La funzione $f(x)$ si avvicina \underline{arbitrariamente}
ad L, a patto che $x$ sia \underline{abbastanza} vicino a $x_0$.

\begin{equation}
    \forall \varepsilon > 0, \exists \delta > 0 :
    \forall n > \delta :
    \forall x \in ] x_0 - \delta, x_0 + \delta [ \backslash \{x_0\}
    \implies f(x) \in ] L - \varepsilon, L + \varepsilon [
\end{equation}

\vspace{5mm}
$\varepsilon$ = intorno al limite

$\delta$ = intorno al valore

L = il limite

\subsection{Serie}

\textbf{DEF}: Sia $(a_n)_n$ una successione di
numeri reali. Definiamo la successione $(s_n)_n$ associata:

\vspace{5mm}
$s_0 = a_0$

\vspace{5mm}
$s_1 = a_0 + a_1$

\vspace{5mm}
$s_2 = a_0 + a_1 + a_2$

\indent
\vdots

\vspace{5mm}
\noindent
La successione delle somme parziali si chiama serie e si indica:

\begin{equation}
    \sum_{n = 0}^{+\infty} a_n \text{ oppure } \sum_{}^{} a_n
\end{equation}

\noindent
I numeri $a_n$ si dicono \underline{termini della serie}.

\noindent
Una serie può essere:

\begin{itemize}
    \item Convergente: $\exists \lim_{n \to \infty} s_n = A \in \mathbb{R}$
    \item Divergente:
    $\exists \lim_{n \to \infty} s_n =
    \begin{cases}
            - \infty \\
            + \infty
    \end{cases}
    $
    \item Indeterminata: $\nexists  \lim_{n \to \infty} s_n$
\end{itemize}

\vspace{5mm}
\noindent
\textbf{DEF}: Data una serie $\sum_{}^{} a_n$, se cambio un
numero finito di termini, la serie generata manterrà
il carattere dell'originale.

\subsubsection{Serie geometrica}

$q \in \mathbb{R}$

\begin{equation}
    R_n := q^n \Rightarrow \sum_{n = 0}^{+\infty} q^n
\end{equation}

\vspace{5mm}
\noindent
La serie geometrica converge soltanto se $|q| < 1$

\begin{equation}
    \sum_{n}^{} q^n = \frac{1}{1 - q}
\end{equation}

\subsubsection{Convergenza di una serie}

Una serie converge se e soltanto se $a_n \to 0$.
La condizione precedente è necessaria ma \underline{NON} sufficiente
per affermare la convergena di una serie.

\subsubsection{Criterio del confronto}

$(a_n)_n, (b_n)_n$

\begin{equation}
    \forall n : 0 \leq a_n \leq b_n
\end{equation}

\begin{itemize}
    \item se $b_n$ converge, anche $a_n$ converge.
    \item se $a_n$ diverge, anche $b_n$ diverge.
\end{itemize}

\subsubsection{Criterio di condensazione}

$\sum_{n}^{} a_n$ converge se e soltanto se $\sum_{k = 0}^{+\infty} 2^k a_2k$ converge.

\subsubsection{Criterio del valore assoluto}

Data una serie $\sum_{n}^{} a_n$, se la serie $\sum_{n}^{} |a_n|$ converge, anche la prima converge.

\begin{equation}
    |\sum_{n}^{} a_n| = \sum_{n}^{} |a_n|
\end{equation}

\noindent
Una serie che viene dichiarata convergente usando questo
criterio è detta assolutamente convergente.

\subsubsection{Criterio della radice}

$(a_n)_n$

\vspace{5mm}
\noindent
se $\exists \lim_{n \to \infty} \sqrt[n]{|a_n|} = L$, allora:

\begin{itemize}
    \item se $L < 1$ allora $\sum_{}^{} a_n$ converge
    \item se $L > 1$ allora $\sum_{}^{} a_n$ \underline{NON} converge
    \item se $L = 1$ allora non si può concludere nulla
\end{itemize}

\subsubsection{Criterio del rapporto}

$(a_n)_n, a_n \neq 0$

\vspace{5mm}
\noindent
se $\exists \lim_{n \to \infty} |\frac{a_n + 1}{a_n}| = L$, allora:

\begin{itemize}
    \item se $L < 1$ allora $\sum_{}^{} a_n$ converge
    \item se $L > 1$ allora $\sum_{}^{} a_n$ \underline{NON} converge
    \item se $L = 1$ allora non si può concludere nulla
\end{itemize}

\subsubsection{Criterio del confronto asintotico}

$(a_n)_n, \forall n : a_n \geq 0$

\noindent
$(b_n)_n, \forall n : b_n > 0$

\vspace{5mm}
\noindent
se $\lim_{n \to \infty} \frac{a_n}{b_n} = L > 0$, allora:

\begin{itemize}
    \item $\sum_{n}^{} a_n < +\infty \Leftrightarrow \sum_{n}^{} b_n < +\infty$, entrambe convergono
    \item $\sum_{n}^{} a_n = +\infty \Leftrightarrow \sum_{n}^{} b_n = +\infty$, entrambe divergono
\end{itemize}

\vspace{5mm}
\noindent
$(a_n)_n, \forall n : a_n \geq 0$

\noindent
$p = n * a_n$

\vspace{5mm}
\noindent
se $\lim_{n \to +\infty} \frac{a_n}{\frac{1}{n^p}} = L > 0$, allora
$\sum_{n}^{} a_n$ converge se $\sum_{n}^{} \frac{1}{n^p}$ converge; possibile soltanto se $p > 1$.

\subsubsection{Criterio di Leibniz}

se $(a_j)_j : a_j > 0 : a_j \geq a_j + 1 : \lim_{j \to \infty} a_j = 0$

\noindent
allora $\sum_{j = 1}^{+\infty} (-1)^j * a_j$ converge.

\end{document}